\section{Start her}
\label{sec:start}

\subsection{Brug under eksamen}

Denne manual er organiseret efter det, du ser i opgaven: differentialligning, blokdiagram,
Bode/Nyquist, performancekrav eller controllerform. Start i metodefinderen, gå til den
angivne opskrift, og skriv dine antagelser før beregningen.

\begin{enumerate}
\item Skriv først præcis hvad opgaven spørger efter: ratio som \(Y/R\), fejl som
      \(E/R\), et gain \(K_P\), en margin, en polplacering eller et tidsresponsmål.
      Marker derefter hvilke signaler der er uafhængige input, og hvilket fortegn hvert
      summationspunkt viser.
\item Klassificer opgaven med tabel~\ref{tab:metodefinder}: vælg rækken ud fra det
      første konkrete tegn i opgaven, fx Bode-plot, Nyquistkurve, blokdiagram,
      differentialligning eller controllerform. Hvis to rækker passer, start med den
      række der bestemmer signalvejen; beregningsformler bruges bagefter.
\item Udfør den manuelle opskrift punkt for punkt. Når et trin henviser til en formel,
      indsæt de kendte størrelser i den nummererede ligning og skriv mellemregningen,
      før du bruger Python som kontrol.
\item Kontrollér stabilitet før DC-gain eller final value theorem bruges: find
      closed-loop-polerne fra den nævner du faktisk bruger, og kræv at alle poler
      ligger strengt i \(\LHP\), medmindre opskriften eksplicit bruger Nyquist.
\item Kontrollér svaret ved fire simple tests: fortegn skal passe med signalretningen,
      enheder skal passe med den efterspurgte størrelse, intervalkrav som
      \(0<\alpha<1\) eller \(\beta>1\) skal være opfyldt, og resultatet skal give mening
      i en grænseværdi som \(K_P=0\), \(\omega\to0\) eller \(\omega\to\infty\).
\end{enumerate}

\subsection{Dokumenteret ramme og kildestatus}

Kurset er \emph{34721/34722 Linear Control Design 1}. Forelæsning 1 angiver en
4-timers skriftlig eksamen suppleret af en assignment-rapport. Det tilgængelige
F21-materiale er multiple choice med løsninger til spørgsmål 1--20. Regler for
hjælpemidler og brug af Python under eksamen fremgår ikke af kilderne; afklar dem
før faktisk eksamensbrug.

\subsection{Tre regler der forhindrer flest fejl}

\begin{itemize}
\item Et plot eller et blokdiagram skal fortolkes manuelt: retning, fortegn og signalplacering
      må ikke antages af et script. Skriv derfor altid en lille signalrelation eller en
      aflæsning, før du regner videre.
\item dB er ikke en lineær gain: brug \eqref{eq:bode-db-conversion} før et gain ganges
      på eller sammenlignes med en margin.
\item Slutværdi kræver stabilt lukket loop; brug poltesten i \S\ref{sec:model} eller
      Nyquist-testen \eqref{eq:nyquist-count} før final value theorem i
      \eqref{eq:steady-state-error}.
\end{itemize}
